% TEMPLATE-MILC-v3.tex - plantilla maestra UNICA de las guias MILC v3.
%
% La usan las tres familias de guias, cada una con su driver:
%   · Grados 6-11  -> scripts/build-guias-g11.py
%   · Bebras       -> scripts/build-guias-bebras.py
%   · Semillero    -> scripts/build-guias-semillero.py
%
% Antes habia tres copias casi identicas (~400 lineas iguales, 6 distintas):
% cada mejora habia que aplicarla tres veces, y en la practica no se hacia
% (el motor de recursos vivia solo en dos de ellas). Lo que varia entre
% familias esta parametrizado en 6 placeholders que cada driver llena
% (se nombran aqui SIN su sintaxis real: el builder reemplaza por texto plano,
% tambien dentro de los comentarios, y un valor multilinea rompe el preambulo):
%
%   HEADER_IZQ        encabezado izquierdo de pagina
%   HEADER_DER        encabezado derecho de pagina
%   PORTADA_ETIQUETA  rotulo sobre el numero grande (GRADO/MOMENTO/MODULO)
%   PORTADA_SERIE     linea de serie de la portada
%   PORTADA_ID        identificador de la guia en la portada
%   PORTADA_CIRCULO   el circulito de grado (°), o vacio si no aplica
%
% El resto de claves las documenta content/guias/_SCHEMA.md. El script valida
% que no queden placeholders sin reemplazar antes de escribir el archivo final.

\documentclass[11pt,letterpaper]{article}
\usepackage[margin=1.75cm]{geometry}
\usepackage[table]{xcolor}
\usepackage{fontspec}
\usepackage[spanish]{babel}
\usepackage{tikz}
% Librerías TikZ para los diagramas de las guías (bloque `recursos:`):
%  · babel  → babel en español vuelve activos caracteres como «>», y TikZ los usa
%             en sus opciones (>=stealth, ->). Sin esta librería, cualquier
%             diagrama con flechas revienta con «\language@active@arg> has an
%             extra }». Debe ir ANTES de que se dibuje nada.
%  · positioning        → sintaxis «below left=2pt and 3pt».
%  · decorations.pathreplacing → llaves ({brace}) para señalar rangos.
\usetikzlibrary{babel,positioning,decorations.pathreplacing}
\usepackage{graphicx}
% Circuitos: el trabajo de electrónica se hace hoy con micro:bit y sus sensores
% integrados (MakeCode, sin cableado), así que no se necesitan esquemáticos.
% Si algún día entran montajes con protoboard, basta descomentar esta línea:
% los diagramas .tex/.tikz declarados en `recursos:` ya se incrustan solos.
% \usepackage{circuitikz}
\usepackage[most]{tcolorbox}
\usepackage{tabularx}
\usepackage{array}
\usepackage{fancyhdr}
\usepackage{titlesec}
\usepackage{ragged2e}
\usepackage{enumitem}
\usepackage{microtype}

% Helvetica Neue no trae la flecha U+2192, asi que cada «→» escrito literalmente
% en el YAML se compilaba a NADA, en silencio (462 flechas en 61 guias: rutas del
% tipo «paleta → tipografia → lockup» salian rotas). Mapeamos el caracter a su
% version matematica, que si tiene glifo. Asi el autor puede seguir escribiendo
% «→» comodamente en el YAML.
\usepackage{newunicodechar}
\newunicodechar{→}{\ensuremath{\rightarrow}}

\setmainfont{Helvetica Neue}
\setsansfont{Helvetica Neue}
\setlength{\parindent}{0pt}
\setlength{\parskip}{6pt}
\hyphenpenalty=200
\exhyphenpenalty=200
\pretolerance=400
\tolerance=2000
\setlength{\emergencystretch}{1em}
\renewcommand{\arraystretch}{1.25}
\setlist{leftmargin=5mm,itemsep=1mm,topsep=1mm}
\newcolumntype{Y}{>{\RaggedRight\arraybackslash}X}

\definecolor{milcMagenta}{HTML}{E6007E}
\definecolor{milcVino}{HTML}{5A0038}
\definecolor{milcTurquesa}{HTML}{0093A5}
\definecolor{milcVerde}{HTML}{58A923}
\definecolor{milcMostaza}{HTML}{E5B400}
\definecolor{milcOcre}{HTML}{B86600}
\definecolor{milcNegro}{HTML}{241816}
\definecolor{milcGris}{HTML}{F7F7F4}
\definecolor{milcLinea}{HTML}{E8E2DD}
\definecolor{milcGradePrimary}{HTML}{0066FF}
\definecolor{milcGradeDark}{HTML}{003D99}
\definecolor{milcGradeSoft}{HTML}{D6E8FF}
\definecolor{milcGradeAccent}{HTML}{CFE7FF}
\definecolor{milcGradeText}{HTML}{FFFFFF}
\color{milcNegro}

\pagestyle{fancy}
\fancyhf{}
\lhead{\small Tecnología e Informática · Grado 11}
\rhead{\small Guía 11 · MILC v3}
\cfoot{\small\thepage}
\renewcommand{\headrulewidth}{0pt}

\newtcolorbox{titlebox}[1]{
  enhanced,colback=#1,colframe=#1,arc=7mm,boxrule=0pt,
  left=7mm,right=7mm,top=3mm,bottom=3mm,
  before skip=8pt,after skip=8pt,
  fontupper=\sffamily\bfseries\Large\color{white}
}

\newtcolorbox{softbox}[3]{
  enhanced,colback=#2,colframe=#1,borderline west={4pt}{0pt}{#1},
  arc=2mm,boxrule=.4pt,left=4mm,right=4mm,top=3mm,bottom=3mm,
  before skip=6pt,after skip=8pt,title={#3},coltitle=white,
  colbacktitle=#1,fonttitle=\sffamily\bfseries,fontupper=\RaggedRight,
  attach boxed title to top left={xshift=3mm,yshift=-2mm},
  boxed title style={arc=2mm, boxrule=0pt}
}

\newtcolorbox{drawbox}[2]{
  enhanced,colback=white,colframe=#1,arc=4mm,boxrule=1pt,
  left=5mm,right=5mm,top=4mm,bottom=4mm,before skip=8pt,after skip=8pt,
  title={#2},coltitle=white,colbacktitle=#1,fonttitle=\sffamily\bfseries,
  fontupper=\RaggedRight,
  attach boxed title to top left={xshift=4mm,yshift=-2mm},
  boxed title style={arc=3mm, boxrule=0pt}
}

\newtcolorbox{infoband}[1]{
  enhanced,colback=#1!8,colframe=#1!50,arc=2mm,boxrule=.5pt,
  left=4mm,right=4mm,top=2mm,bottom=2mm,before skip=4pt,after skip=4pt,
  fontupper=\small\RaggedRight
}

% Puente narrativo entre secciones (contrato editorial Conéctate).
% Da continuidad: "Acabas de X. Ahora vas a Y porque Z."
% Visualmente: banda izquierda en color del grado, texto en itálica pequeña.
\newtcolorbox{puentebox}{
  enhanced,colback=milcGradeSoft,colframe=milcGradeDark,
  borderline west={3pt}{0pt}{milcGradeDark},
  arc=0mm,boxrule=0pt,left=5mm,right=4mm,top=2.5mm,bottom=2.5mm,
  before skip=4pt,after skip=8pt,
  fontupper=\itshape\small\color{milcNegro}\RaggedRight
}
% La flecha va en modo matematico a proposito: Helvetica Neue NO tiene el glifo
% U+2192, asi que \textrightarrow se compilaba a NADA («Missing character: There
% is no → in font Helvetica Neue») y el puente salia sin flecha. En modo
% matematico la flecha viene de la fuente matematica y si se dibuja.
\newcommand{\puente}[1]{%
  \begin{puentebox}%
    \textcolor{milcGradeDark}{$\rightarrow$}\hspace{2mm}#1%
  \end{puentebox}%
}

\newcommand{\linead}[1]{\noindent\color{milcLinea}\rule{#1}{0.5pt}\color{milcNegro}}
\newcommand{\respuesta}[1]{\par\vspace{2mm}\foreach \n in {1,...,#1}{\noindent\color{milcLinea}\rule{\linewidth}{0.5pt}\par\vspace{5mm}}\color{milcNegro}}
\newcommand{\checkbox}{\textcolor{milcGradeDark}{\fbox{\phantom{x}}}\hspace{5pt}}

\newcommand{\phasecard}[4]{
\begin{tcolorbox}[enhanced,colback=#3,colframe=#2,arc=3mm,boxrule=.5pt,height=3.05cm,valign=center,halign=center,left=1.5mm,right=1.5mm,top=2mm,bottom=2mm]
{\sffamily\bfseries\fontsize{22}{24}\selectfont\textcolor{#2}{#1}}\par
{\sffamily\bfseries\small #4}
\end{tcolorbox}
}

% ── Recursos visuales de la guía ──────────────────────────────────────
% Declarados en el bloque `recursos:` del YAML y emitidos por el builder.
% \guiaFigura   → imagen rasterizada o PDF (foto, carta celeste, montaje).
% \guiaDiagrama → fuente TikZ/circuitikz (.tex) incrustada como vector:
%                 es la vía para circuitos y esquemáticos editables.
% #1 = ruta relativa al .tex · #2 = pie (puede ir vacío)
\newcommand{\guiaFigura}[2]{%
\begin{center}
\includegraphics[width=0.88\linewidth,height=0.38\textheight,keepaspectratio]{#1}\\[1.5mm]
{\footnotesize\itshape\textcolor{milcNegro!75}{#2}}
\end{center}
}

\newcommand{\guiaDiagrama}[2]{%
\begin{center}
\input{#1}\\[1.5mm]
{\footnotesize\itshape\textcolor{milcNegro!75}{#2}}
\end{center}
}

\begin{document}
\thispagestyle{empty}
\begin{tikzpicture}[remember picture,overlay]
  \fill[white] (current page.south west) rectangle (current page.north east);
  \path[fill=milcGradePrimary,rounded corners=16mm]
    ([xshift=.55cm,yshift=-.55cm]current page.north west) rectangle
    ([xshift=-.55cm,yshift=3.65cm]current page.south east);

  \node[anchor=north west,text=milcGradeText] at ([xshift=2.0cm,yshift=-1.85cm]current page.north west)
    {\sffamily\bfseries\fontsize{14}{16}\selectfont GRADO};
  \node[anchor=north west,text=milcGradeText,opacity=.82] at ([xshift=2.0cm,yshift=-2.75cm]current page.north west)
    {\sffamily\bfseries\fontsize{9}{11}\selectfont SERIE GUÍAS · TECNOLOGÍA E INFORMÁTICA};

  \begin{scope}[shift={([xshift=-3.05cm,yshift=-2.55cm]current page.north east)}]
    \fill[white,opacity=.80] (-.62,-.52) rectangle (.62,.52);
    \draw[milcGradeText,line width=1pt] (-.62,-.52) rectangle (.62,.52);
    \fill[milcGradeDark] (-.42,-.38) rectangle (-.22,.06);
    \fill[milcGradeAccent] (-.10,-.38) rectangle (.10,.24);
    \fill[milcGradePrimary!60!black] (.22,-.38) rectangle (.42,.38);
  \end{scope}

  \node[anchor=west,text=milcGradeText] at ([xshift=1.9cm,yshift=-12.85cm]current page.north west)
    {\sffamily\bfseries\fontsize{124}{124}\selectfont 11};
  % El circulito de grado (°) solo aplica a las guias de grado («11°»); en
  % Bebras y Semillero el numero grande es un momento/modulo, no un grado.
  \draw[milcGradeText,line width=1.7pt]
    ([xshift=4.52cm,yshift=-11.68cm]current page.north west) circle (.13cm);

  \node[anchor=west,text=milcGradeText,text width=8.7cm,align=left] at ([xshift=10.35cm,yshift=-11.6cm]current page.north west)
    {
     {\sffamily\bfseries\fontsize{19}{22}\selectfont Mapeo de procesos ---\\ver dónde\\se pierde el tiempo}\\[4mm]
     {\sffamily\bfseries\fontsize{23}{26}\selectfont Guía 11-11-TIC}\\[2mm]
     {\sffamily\fontsize{13}{16}\selectfont Período 2 · Automatización y procesos}\\[5mm]
     {\sffamily\bfseries\fontsize{11}{13}\selectfont Institución Educativa Sor María Juliana}\\[1mm]
     {\sffamily\fontsize{10}{12}\selectfont Autor: Dr. Álvaro Cárdenas Orozco}
    };

  \path[fill=milcGradeSoft,rounded corners=7mm]
    ([xshift=1.35cm,yshift=.85cm]current page.south west) rectangle
    ([xshift=-1.35cm,yshift=4.25cm]current page.south east);
  \draw[milcGradeDark,line width=.65pt,rounded corners=7mm]
    ([xshift=1.35cm,yshift=.85cm]current page.south west) rectangle
    ([xshift=-1.35cm,yshift=4.25cm]current page.south east);
  \node[anchor=center,text=milcNegro,text width=17.0cm,align=center] at ([yshift=2.55cm]current page.south)
    {\sffamily\fontsize{10}{12}\selectfont Metodología MILC · Apertura ancestral · Pensamiento computacional · Triángulo Dussel-Estoicismo-Floridi\\
    \textcolor{milcGradeDark}{\bfseries Producto:} Mapa visual de un proceso real (escuela, casa, microempresa cercana) con todos los pasos, responsables, tiempos estimados y al menos 2 cuellos de botella marcados con justificación de por qué lo son};
\end{tikzpicture}
\mbox{}
\newpage

\section*{Apertura: Mapeo de procesos --- ver dónde se pierde el tiempo}

\begin{softbox}{milcGradeDark}{milcGradeSoft}{Saber ancestral}
El \textbf{maestro carpintero} de Cartago no necesita reloj para saber dónde se traba el taller: lo ve. Sabe que el aserrado se demora porque la sierra está mal afilada, que el lijado se acumula porque solo un aprendiz lo hace, que el barnizado se traba porque no hay ventilación. La \textbf{cocinera mayor} de una fonda de mercado ve el cuello de botella sin medirlo: \emph{``el sancocho se demora porque solo hay una estufa''}. El \textbf{maestro del oficio} no necesita lenguaje técnico para ver el flujo: lleva en la mirada el saber acumulado de cien repeticiones. Esa mirada es el primer acto de toda automatización: \emph{antes de cambiar el proceso, hay que verlo}.
\end{softbox}

\begin{softbox}{milcTurquesa}{milcGris}{Saber contemporáneo}
El \textbf{mapeo de procesos} es la representación visual de cómo se hace algo, paso a paso, con responsables y tiempos. Es la traducción técnica de la mirada del maestro al lenguaje del análisis. Un mapa de proceso responde a 4 preguntas: \emph{¿qué pasos hay?}, \emph{¿quién hace cada uno?}, \emph{¿cuánto tarda cada uno?}, \emph{¿dónde se acumula el trabajo?}. El cuarto punto --- el \textbf{cuello de botella} --- es donde se pierde la mayoría del tiempo total. La regla del oficio: \emph{automatizar antes de mapear es perder el tiempo en lo que no importa}. Sin mapa, terminas comprando un robot para el paso que tardaba 2 minutos en lugar de arreglar el que tardaba 20.
\end{softbox}

\begin{softbox}{milcMostaza}{milcGris}{Pregunta-puente}
\textit{¿Cómo sabía la cocinera mayor, sin cronómetro, que el sancocho era el cuello de botella de la fonda? ¿Y qué pierde un emprendedor novato cuando compra software de automatización sin haber mapeado primero su proceso?}
\end{softbox}

\begin{softbox}{milcVerde}{milcGris}{Saber y saber hacer}
Al terminar podrás: (1) \textbf{identificar} los pasos, responsables y tiempos de un proceso real cercano; (2) \textbf{analizar} dónde se acumula el trabajo y por qué (cuello de botella); (3) \textbf{explicar} con tus palabras qué cambia si se interviene en el cuello de botella vs en otro paso; (4) \textbf{evaluar} qué 2 cuellos de botella son los más graves y qué tipo de intervención (automatizar, redistribuir, eliminar) tendría más impacto.
\end{softbox}



\small
\begin{tabularx}{\textwidth}{>{\bfseries}p{3.0cm}Y}
\rowcolor{milcGradePrimary!12}Autor & Dr. Álvaro Cárdenas Orozco\\
DBA & Diseño y mejora de procesos digitales con criterio ético (MEN, T\&I)\\
Referentes & Pensamiento computacional · Lean / Kaizen · Floridi (ética informacional)\\
Producto final & Mapa visual de un proceso real (escuela, casa, microempresa cercana) con todos los pasos, responsables, tiempos estimados y al menos 2 cuellos de botella marcados con justificación de por qué lo son\\
\end{tabularx}
\normalsize

\puente{Antes de empezar, mira el viaje completo. En el periodo 1 construiste tu presencia digital personal/marca. En este periodo 2 vas a mirar adentro de un proceso real: cómo funciona, dónde falla, qué se puede automatizar. Hoy abres ese trabajo con la mirada del maestro: \emph{ver antes de cambiar}.}

\begin{titlebox}{milcGradeDark}
Ruta MILC: así vamos a aprender
\end{titlebox}
\begin{tabularx}{\textwidth}{>{\centering\arraybackslash}X>{\centering\arraybackslash}X>{\centering\arraybackslash}X>{\centering\arraybackslash}X}
\phasecard{1}{milcVerde}{milcGris}{Escucha\\[-1mm]\small Observo, escucho y pregunto.} &
\phasecard{2}{milcTurquesa}{milcGris}{Sistematización\\[-1mm]\small Pensamiento computacional.} &
\phasecard{3}{milcMagenta}{milcGris}{Praxis\\[-1mm]\small Construyo y aplico.} &
\phasecard{4}{milcVino}{milcGris}{Evaluación\\[-1mm]\small Reflexión liberadora.}
\end{tabularx}


\puente{Antes de teorizar, vas a observar un proceso real de tu día (preparar el desayuno, llegar al colegio, registrar la nota de un examen) y a anotar todos los pasos sin filtrar. El reconocimiento del flujo en bruto es la materia prima del mapa.}

\begin{titlebox}{milcVerde}
Fase 1 · Escucha
\end{titlebox}

\begin{softbox}{milcVerde}{milcGris}{Escena para escuchar}
\textbf{Actividad 1 · IDENTIFICA --- Caza el cuello de botella} (15 min · individual). Elige un proceso real de tu día (preparar el desayuno, llegar al colegio, registrar la nota de un examen en la plataforma, despachar un pedido en una tienda del barrio). Sin filtro, anota TODOS los pasos en orden, desde el primero hasta el último. Junto a cada paso, escribe quién lo hace y cuánto tarda. Al final, marca con 🐌 el paso que \emph{intuyes} que es el cuello de botella, sin medirlo aún.
\end{softbox}

\checkbox Listé todos los pasos del proceso en orden, sin omitir ninguno.\\
\checkbox Anoté quién hace cada paso y cuánto estimo que dura.\\
\checkbox Marqué con 🐌 mi candidato a cuello de botella e intenté justificar por qué.

\begin{infoband}{milcVerde}


\textbf{Lo que va al cuaderno (Actividad 1 · IDENTIFICA).} Título: «Actividad 1 --- Caza el cuello de botella». \textbf{Formato:} tabla 4 columnas (Paso~|~Responsable~|~Tiempo estimado~|~Notas), 8-15 filas según el proceso. Al final marca 🐌 en el paso sospechoso. \textbf{Sabes que terminaste cuando} puedes contarle a un compañero el proceso completo sin volver a mirar la tabla.
\end{infoband}

\puente{Ya viste el flujo. Ahora vas a entender la \textbf{anatomía} de un mapa de proceso: cómo representar pasos, responsables y tiempos; cómo detectar cuellos de botella; cómo no confundir \emph{``tarda mucho''} con \emph{``es el más importante para optimizar''}.}

\begin{titlebox}{milcTurquesa}
Fase 2 · Sistematización con pensamiento computacional
\end{titlebox}

\begin{softbox}{milcTurquesa}{milcGris}{Cuatro pilares aplicados al tema}
Un mapa de proceso bien hecho no parece complicado, pero su anatomía es estricta. Vamos a descomponerla con los pilares del pensamiento computacional para que no quede al ojo del aprendiz, sino al rigor del oficio.
\end{softbox}

\small
\begin{tabularx}{\textwidth}{>{\bfseries\RaggedRight\arraybackslash}p{3.6cm}Y}
\rowcolor{milcTurquesa!90}\color{white}Pilar & \color{white}\bfseries Aplicación al tema\\
1. Descomposición & El mapa se descompone en 4 elementos visibles: (1) \textbf{Pasos}: cada acción concreta numerada (no \emph{``trabajar''}, sino \emph{``cortar la tabla a 50 cm''}); (2) \textbf{Responsables}: quién hace cada paso (persona, rol, máquina); (3) \textbf{Tiempos}: cuánto tarda cada paso en promedio; (4) \textbf{Flujo}: las flechas que muestran qué pasa después de qué (incluye bucles, esperas, decisiones).\\
2. Reconocimiento de patrones & Todo cuello de botella sigue el patrón \textbf{``el paso lento $\to$ marca el ritmo del todo''}: si un paso tarda 20 min y todos los demás 2 min, la velocidad de \emph{todo el proceso} es la velocidad del paso de 20 min. Aumentar la velocidad de los pasos rápidos no cambia el total. \emph{Optimizar lo que no es cuello de botella es perder el tiempo}. Esa es la regla del oficio que distingue al ingeniero del que solo \emph{``compra herramientas nuevas''}.\\
3. Abstracción & Lo esencial cabe en una pregunta: \emph{¿qué pasaría con todo el proceso si el cuello se resuelve?}. Si la respuesta es \emph{``baja el tiempo total a la mitad''}, el cuello es grave. Si la respuesta es \emph{``ahorro 2 minutos pero el total queda igual''}, el cuello no era el cuello. La phronesis del análisis exige medir el impacto, no la severidad subjetiva.\\
4. Algoritmo conceptual & Algoritmo para encontrar el cuello: (1) suma los tiempos de todos los pasos; (2) calcula qué porcentaje del total ocupa cada paso; (3) identifica el paso que ocupa más de 30\% del total (el primer candidato); (4) verifica si tiene espera, acumulación o requiere intervención manual; (5) compara con el segundo candidato; (6) marca los 2 cuellos más graves.\\
\end{tabularx}
\normalsize

\begin{drawbox}{milcTurquesa}{Anatomía de un mapa de proceso bien hecho}
\textbf{Título:} nombre del proceso en 1 frase. \\[1mm] \textbf{Pasos numerados:} en orden, cada uno con verbo de acción concreto. \\[1mm] \textbf{Responsables:} persona, rol o máquina por paso. \\[1mm] \textbf{Tiempos:} duración estimada por paso. \\[1mm] \textbf{Flujo:} flechas + decisiones + bucles si los hay. \\[1mm] \textbf{Cuello de botella:} marcado con color, con justificación de cuánto ocupa del total.
\end{drawbox}

\begin{softbox}{milcMostaza}{milcGris}{Errores comunes y su corrección}
(1) Confundir \emph{``el paso que más me molesta''} con \emph{``el cuello de botella''}: lo que te molesta puede ser irrelevante para el tiempo total. (2) Omitir las esperas: si entre el paso 3 y el 4 hay 1 hora de espera porque \emph{``el papel debe secar''}, esa hora cuenta. (3) Confundir responsable con beneficiario: el responsable es quien \emph{ejecuta} el paso, no quien recibe el resultado. (4) Mapear con tiempos imaginarios: si no mediste, declara \emph{``estimado''}; si mediste, declara la fuente. (5) Marcar 5 cuellos: si todo es cuello, nada es cuello.
\end{softbox}

\begin{infoband}{milcTurquesa}


\textbf{Lo que va al cuaderno (Actividad 2 · ANALIZA).} Título: «Actividad 2 --- Anatomía del cuello de botella». \textbf{Formato:} ficha con las 6 partes del mapa (título, pasos, responsables, tiempos, flujo, cuello) aplicadas al proceso de Actividad 1 + 1 cálculo de \% del total que ocupa cada paso. \textbf{Sabes que terminaste cuando} puedes señalar con números cuál es el cuello y cuál es solo un paso \emph{``molesto pero rápido''}.
\end{infoband}

\puente{Tienes el método. Toca aplicarlo: elegir un proceso real cercano, mapear sus pasos, marcar tiempos, identificar responsables, encontrar al menos 2 cuellos de botella y justificarlos. Producto que puedes usar como evidencia ante un emprendedor real.}

\begin{titlebox}{milcMagenta}
Fase 3 · Praxis con pensamiento computacional
\end{titlebox}

\begin{softbox}{milcMagenta}{milcGris}{Construyendo el producto con los cuatro pilares}
\textbf{Actividad 3 · EXPLICA + EVALÚA --- Mapa con 2 cuellos justificados} (30 min · individual). Hora de aplicar. Elige un proceso real cercano (preferiblemente de un negocio del barrio, una microempresa familiar o un proceso escolar) y produce el mapa completo con 2 cuellos de botella identificados y justificados.
\end{softbox}

\small
\begin{tabularx}{\textwidth}{>{\bfseries\RaggedRight\arraybackslash}p{3.6cm}Y}
\rowcolor{milcMagenta!90}\color{white}Pilar & \color{white}\bfseries Aplicación al producto\\
Descomposición & Tu mapa se descompone en 5 piezas: (1) título del proceso + responsable general; (2) tabla de pasos con tiempo y responsable; (3) representación visual (diagrama a mano alzada o digital); (4) cálculo de \% del total por paso; (5) 2 cuellos de botella marcados con justificación.\\
Patrones & Sigue el patrón \textbf{``observa primero, mapea después''}: si el proceso no es tuyo, ve y observa al menos una corrida completa antes de mapear. Mapear de memoria o por suposición produce mapas inventados. La observación directa toma 30 min pero ahorra 3 horas de mal análisis.\\
Abstracción & Cada justificación expresa una sola idea: \emph{``Este es cuello porque ocupa X\% del tiempo total Y por la espera de Z''}. Si tu justificación tiene 3 razones, parte el cuello en 2: probablemente son 2 problemas distintos juntos.\\
Algoritmo de armado & Algoritmo: (1) elige el proceso; (2) observa una corrida completa (o entrevista a quien lo hace); (3) lista pasos con tiempos; (4) dibuja el flujo; (5) calcula \%; (6) identifica los 2 cuellos; (7) justifica cada uno en 1 frase con número + razón.\\
\end{tabularx}
\normalsize

\begin{drawbox}{milcMagenta}{Checklist antes de entregar el mapa}
\checkbox Proceso observado al menos 1 vez completo (no mapeado de memoria). \\[1mm] \checkbox Todos los pasos numerados con responsable y tiempo estimado. \\[1mm] \checkbox Diagrama visual con flujo y decisiones marcadas. \\[1mm] \checkbox Cálculo de \% del tiempo total por paso. \\[1mm] \checkbox 2 cuellos marcados con justificación en 1 frase cada uno. \\[1mm] \checkbox Un emprendedor podría decidir dónde intervenir sin pedirte explicación.
\end{drawbox}

\begin{softbox}{milcMagenta}{milcGris}{Plantilla de trabajo}
\textbf{Proceso.} \textit{Nombre + responsable general + dónde se ejecuta.} \\[1mm] \textbf{Pasos (tabla).} \textit{Cada uno con número, descripción, responsable, tiempo, \% del total.} \\[1mm] \textbf{Diagrama.} \textit{Boceto del flujo (rectángulos para pasos, rombos para decisiones).} \\[1mm] \textbf{Cuello 1.} \textit{Paso N + justificación numérica.} \\[1mm] \textbf{Cuello 2.} \textit{Paso M + justificación numérica.}
\end{softbox}

\begin{infoband}{milcMagenta}


\textbf{Lo que va al cuaderno (Actividad 3 · EXPLICA + EVALÚA).} Título: «Actividad 3 --- Mapa de proceso con cuellos de botella». \textbf{Formato:} 1 página con título + tabla de pasos + boceto del diagrama + 2 cuellos justificados. \textbf{Extensión:} 1 página A4 o 2 páginas de cuaderno. \textbf{Sabes que terminaste cuando} podrías presentar el mapa al responsable real del proceso y proponer una conversación de mejora basada en evidencia.
\end{infoband}

\puente{Lo que sigue resume \emph{qué} entregas hoy: 1 mapa visual + 2 cuellos de botella justificados. El criterio principal: que un emprendedor del barrio, al ver tu mapa, pueda \textbf{decidir} dónde intervenir primero sin pedirte explicación.}

\begin{titlebox}{milcMostaza}
Producto de la sesión
\end{titlebox}
\textbf{1 mapa de proceso real con 2 cuellos de botella justificados}.
\begin{softbox}{milcMostaza}{milcGris}{Criterios de calidad}
(1) Proceso real elegido y observado (no inventado de memoria). (2) Tabla de pasos con responsable y tiempo por paso. (3) Diagrama visual con flujo y decisiones. (4) Cálculo de \% del tiempo total por paso. (5) 2 cuellos de botella marcados con justificación numérica.
\end{softbox}

\puente{Antes de cerrar, mira el mapeo desde las cinco dimensiones humanas. Mapear bien es respeto al oficio; mapear mal es desprecio por la sabiduría acumulada del que ya hace el trabajo.}

\begin{titlebox}{milcVino}
Fase 4 · Evaluación liberadora — cinco dimensiones
\end{titlebox}

\begin{softbox}{milcVino}{milcGris}{Sentido de la evaluación}
Evaluar no es solo revisar una nota. Es reconocer qué aprendí, qué cambió en mí, qué emociones manejé, cómo me relaciono con la ciudad y la comunidad, y qué le dejaré a quien viene detrás.
\end{softbox}

\textbf{Mi reflexión liberadora en cinco dimensiones.} Completa cada frase iniciada:

\small
\begin{tabularx}{\textwidth}{>{\bfseries\RaggedRight\arraybackslash}p{3.4cm}Y>{\RaggedRight\arraybackslash}p{4.6cm}}
\rowcolor{milcVino}\color{white}Dimensión & \color{white}\bfseries Mi respuesta (frame starter) & \color{white}\bfseries Algo que descubrí\\
1. Personal & Lo que más cambió en mí fue \linead{6.5cm} & \linead{4.2cm}\\
2. Emocional & La emoción más fuerte fue \linead{2.5cm} y la regulé \linead{3cm} & \linead{4.2cm}\\
3. Ciudadana & En democracia esto importa porque \linead{6cm} & \linead{4.2cm}\\
4. Local (Cartago) & En mi barrio sería útil para \linead{6.5cm} & \linead{4.2cm}\\
5. Intergeneracional & A mi abuela/abuelo le diría \linead{6.5cm} & \linead{4.2cm}\\
\end{tabularx}
\normalsize

\begin{infoband}{milcVino}
\textbf{Para la próxima sustentación.} Vuelve a esta tabla en seis meses. Verás que las respuestas cambian — no porque lo aprendido cambie, sino porque tú cambias. La evaluación liberadora es una conversación contigo a través del tiempo.
\end{infoband}


\puente{Y para terminar, tres voces sobre ver el proceso. Dussel sobre los procesos que invisibilizan a quien sostiene el trabajo, Marco Aurelio sobre la disciplina de observar antes de intervenir, Floridi sobre el mapa como infraestructura cognitiva del oficio contemporáneo.}

\begin{titlebox}{milcGradeDark}
Triángulo de pensamiento · cierre filosófico
\end{titlebox}

\begin{softbox}{milcVino}{milcGris}{Dussel · lente del nosotros}
\textit{``Todo mapa de proceso decide quién aparece como protagonista del trabajo y quién queda invisible.''}

Cuando mapeas un proceso de una panadería y nombras \emph{``el panadero''} pero no nombras a \emph{``la señora que limpia el horno cada noche para que el panadero pueda trabajar mañana''}, tu mapa invisibiliza el trabajo de cuidado. La phronesis del mapeo exige nombrar a todos los responsables, incluyendo los del trabajo invisible (limpieza, mantenimiento, abastecimiento, transporte). Un mapa que solo nombra al protagonista repite la estructura social que lo hace protagonista.

\textbf{En mi mapa, ¿quién aparece nombrado y quién quedó como \emph{``ya estaba listo''} sin responsable? ¿Qué cambia si nombro al trabajo invisible que sostiene el proceso?}
\end{softbox}
\linead{\linewidth}\\[1mm]
\linead{\linewidth}

\begin{softbox}{milcMostaza}{milcGris}{Estoicismo · Marco Aurelio · cuidado interior}
\textit{``Antes de cambiar el flujo del trabajo, observa el flujo del trabajo.''}

Es tentador llegar con propuestas de automatización antes de haber visto el proceso. La sabiduría estoica recomienda lo contrario: \emph{quien interviene sin observar daña más de lo que arregla}. El maestro del oficio sabe que el aprendiz que llega con soluciones está más perdido que el que llega con preguntas. Mapear es la disciplina de \emph{observar antes de intervenir}.

\textbf{¿Llegué al mapa con la respuesta ya decidida o llegué con preguntas? ¿Qué descubrí en la observación que no habría podido inventar desde el escritorio?}
\end{softbox}
\linead{\linewidth}\\[1mm]
\linead{\linewidth}

\begin{softbox}{milcTurquesa}{milcGris}{Floridi · lente de la infoesfera}
\textit{``El mapa de procesos es la infraestructura cognitiva con que toda organización contemporánea se entiende a sí misma.''}

Cualquier organización pública o privada que automatiza, optimiza o se transforma digitalmente parte de un mapa de sus procesos. Sin mapa, las decisiones tecnológicas son adivinanzas caras. Aprender a mapear procesos a los 16 años es alfabetización profesional anticipada: el lenguaje con el que la ingeniería, la administración y la gestión pública entienden las organizaciones.

\textbf{¿Cuántos procesos cotidianos podría mejorar (en mi casa, mi colegio, el negocio familiar) si me tomara 1 hora de mapeo en lugar de quejarme de su lentitud?}
\end{softbox}
\linead{\linewidth}\\[1mm]
\linead{\linewidth}

\begin{titlebox}{milcVino}
Compromiso final
\end{titlebox}

\small
\begin{tabularx}{\textwidth}{>{\RaggedRight\arraybackslash}p{3.0cm}YYY}
\rowcolor{milcVino}\color{white}\bfseries Criterio & \color{white}\bfseries Lo logré & \color{white}\bfseries En proceso & \color{white}\bfseries Necesito apoyo\\
Comprensión del tema & Explico ideas centrales con mis palabras. & Reconozco ideas pero falta precisión. & Necesito repasar conceptos base.\\
Pensamiento computacional & Apliqué los 4 pilares al diseñar mi producto. & Apliqué algunos pilares. & Necesito releer la fase 2.\\
Aplicación y producto & Mi producto cumple los criterios y se puede revisar. & Producto funcional con ajustes pendientes. & Aún en construcción.\\
Reflexión 5 dimensiones & Respondí cada dimensión con honestidad. & Respondí algunas con detalle. & Mis respuestas son muy generales.\\
Triángulo de pensamiento & Conecté las 3 voces con mi trabajo real. & Conecté al menos una. & No relaciono las voces con mi proceso.\\
\end{tabularx}
\normalsize

\textbf{Hoy entendí que:}\\[1mm]
\linead{\linewidth}\\[1mm]
\linead{\linewidth}

\textbf{Mi compromiso para la próxima sesión es:}\\[1mm]
\linead{\linewidth}\\[1mm]
\linead{\linewidth}

\begin{softbox}{milcMostaza}{milcGris}{Cita de despedida · Marco Aurelio}
\textit{``El obstáculo en el camino se vuelve el camino. Lo que estorba al camino se vuelve camino.''}\\[1mm]
Cada error de hoy, bien observado, se convierte en pieza del camino que pisarás mañana.
\end{softbox}

\small
\begin{tabularx}{\textwidth}{>{\bfseries}p{3.6cm}Y}
\rowcolor{milcGradeDark!12}Nombre completo: & \linead{10cm}\\
Fecha: & \linead{4cm} \quad Curso/grupo: \linead{4cm}\\
Mi firma: & \linead{9cm}\\
\end{tabularx}
\normalsize

\begin{infoband}{milcGradeDark}
\textbf{Para la próxima semana.} Mapea un segundo proceso de tu vida cotidiana (estudiar para un examen, organizar un evento, hacer las compras) y comparte el mapa con quien comparta el proceso contigo (madre, padre, hermano, compañero). La phronesis del mapeo no se queda en el cuaderno: se conversa con quien sostiene el trabajo.
\end{infoband}

\end{document}
